\documentclass[11pt,a4paper]{article}

\usepackage[T1]{fontenc}
\usepackage[utf8]{inputenc}
\usepackage[a4paper,margin=2.2cm]{geometry}
\usepackage{microtype}
\usepackage{lmodern}
\usepackage{amsmath,amssymb}
\usepackage{booktabs}
\usepackage{longtable}
\usepackage{array}
\usepackage{tabularx}
\usepackage{xcolor}
\usepackage{listings}
\usepackage{hyperref}
\usepackage{enumitem}
\usepackage{titlesec}
\usepackage{fancyhdr}

\definecolor{codebg}{RGB}{245,245,245}
\definecolor{codeframe}{RGB}{200,200,200}
\definecolor{warnred}{RGB}{180,30,30}
\definecolor{okgreen}{RGB}{30,120,40}
\definecolor{accentblue}{RGB}{30,80,160}

\hypersetup{
  colorlinks=true,
  linkcolor=accentblue,
  urlcolor=accentblue,
  citecolor=accentblue
}

\lstdefinestyle{plain}{
  basicstyle=\ttfamily\footnotesize,
  backgroundcolor=\color{codebg},
  frame=single,
  rulecolor=\color{codeframe},
  breaklines=true,
  breakatwhitespace=false,
  columns=fullflexible,
  keepspaces=true,
  showstringspaces=false,
  xleftmargin=4pt,
  xrightmargin=4pt,
  framexleftmargin=2pt,
  framexrightmargin=2pt
}
\lstdefinestyle{py}{
  style=plain,
  language=Python,
  keywordstyle=\color{accentblue}\bfseries,
  commentstyle=\color{gray}\itshape,
  stringstyle=\color{okgreen}
}
\lstset{style=plain}

\titleformat{\section}{\Large\bfseries\color{accentblue}}{\thesection}{0.6em}{}
\titleformat{\subsection}{\large\bfseries}{\thesubsection}{0.6em}{}

\pagestyle{fancy}
\fancyhf{}
\fancyhead[L]{\small DAMPS--MMHCL Phase~1 -- H2 Diagnostic Positive}
\fancyhead[R]{\small rev44 / Clothing}
\fancyfoot[C]{\thepage}

\title{\textbf{Phase~1 Day-1 Diagnostic Sprint: \\
H2 Diagnostic Positive on Amazon Clothing} \\[0.3em]
\large An In-Depth Analysis of the \texttt{return code 1} Outcome \\
\large and Its Implications for the rev44 Phase~1 Stop-Gate Verdict}
\author{DAMPS--MMHCL Project Team}
\date{May 12, 2026}

\begin{document}
\maketitle

\begin{center}
\fbox{\parbox{0.92\textwidth}{\textbf{Bottom line.} The non-zero return code reported at the end of the
$\sim$12-hour Day-1 sprint is \emph{not} an infrastructure failure. It is the
\textbf{intended diagnostic-positive exit code} for hypothesis \textbf{H2}
defined in the Phase~1 Root-Cause Roadmap, §2.1. Because \texttt{meta\_categories.npy}
is absent on disk, \texttt{utility/load\_data.py} silently fell back to a deterministic
hash and \textbf{every} APC-driven training run, including all 20 runs of the
Phase~1 four-variant sweep in cell~16, was executed on phantom clusters.
The Phase~1 stop-gate verdict is therefore \textbf{confounded by H2}
and must be re-evaluated after the real metadata file is restored.}}
\end{center}

\section{TL;DR}

\begin{itemize}[leftmargin=*]
  \item \textbf{Return code 1 $\neq$ an operational error.} It is precisely the
        \textbf{H2 SUSPECTED} signal from roadmap §2.1 -- the Day-1 script is
        telling you: ``\textit{the metadata file is missing; APC has been
        running on a hash-based fallback from epoch~0}.''
  \item This is a \textbf{positive diagnostic} for H2 among the ten root-cause
        hypotheses in the roadmap, and it \textbf{automatically invalidates a
        portion of the Phase~1 conclusions}, because all 20 multi-seed training
        runs in cell~16 also executed with the same missing-meta configuration --
        meaning APC produced phantom signals in \textbf{every} variant.
\end{itemize}

\section{Interpreting the Warning Banner}

\begin{lstlisting}
[pre-flight] MISSING ...\data\Clothing\meta_categories.npy
return code: 1
[WARN] Non-zero exit code -- review logs above.
\end{lstlisting}

\noindent Cross-referenced against the source of
\texttt{day1\_diagnostic\_sprint.py} (lines 106--113):

\begin{lstlisting}[style=py]
if not exists:
    report["h2_suspected"] = True
    report["warnings"].append(
        "meta_categories.npy MISSING -> load_data.py uses deterministic hash "
        "fallback (phantom APC clusters; roadmap H2)."
    )
\end{lstlisting}

\noindent and lines 305--306:

\begin{lstlisting}[style=py]
if rep["h2_suspected"]:
    exit_code = 1
\end{lstlisting}

\noindent $\Rightarrow$ The exit code \texttt{1} here is \textbf{not a crash}; it is the
``diagnostic-positive'' exit code that the roadmap deliberately uses for
\textbf{fail-fast} signalling. D2, D3 and the APC-off probe \textbf{still run
after D1} (see lines 310, 340 and 370 of the script -- none of those branches
is short-circuited by a \texttt{return}).

\subsection*{Natural consequences}

\begin{itemize}[leftmargin=*]
  \item The $\sim$12-hour wall-clock run almost certainly completed
        D2 (4 values of $K$, 1 run each) + D3 (1 run) + the APC-off probe
        (1 run) = \textbf{6 GPU runs total}.
  \item Only the opening banner and the closing return code were pasted -- the
        body of stdout, which contains the D2/D3/APC results, was truncated
        during the copy.
  \item The full diagnostic conclusion requires reading the per-run log files
        (see the \textbf{Action items} section below).
\end{itemize}

\section{Serious Implication for the Phase~1 Sweep (Cell 16)}

The Phase~1 sweep in cell~16 (20 runs: 5 seeds $\times$ 4 variants) was
executed \textbf{before} cell~15 on the \textbf{same machine} with the
\textbf{same data directory}. Because \texttt{meta\_categories.npy} did not
exist, \texttt{utility/load\_data.py} silently entered its hash-fallback branch
for \textbf{all 20 runs}. The per-variant impact is summarised in
Table~\ref{tab:variant-impact}.

\begin{table}[h]
\centering
\renewcommand{\arraystretch}{1.25}
\begin{tabularx}{\textwidth}{|l|c|c|c|X|}
\hline
\textbf{Variant} & \texttt{damps\_apc} & \textbf{Hash-fallback?} &
\textbf{Measured R@20} & \textbf{Revised verdict} \\
\hline
(a) anchor    & =1 & Yes (APC on, fake clusters) & $0.0795 \pm 0.0025$ &
APC signal contaminated \\
\hline
(b) tau03     & =1 & Yes & $0.0829 \pm 0.0007$ &
APC signal contaminated \\
\hline
(c) avrf\_off & =1 & Yes & $0.0803 \pm 0.0021$ &
APC signal contaminated \\
\hline
(d) combined  & =1 & Yes & $\mathbf{0.0851 \pm 0.0012}$ &
APC signal contaminated \\
\hline
\end{tabularx}
\caption{Per-variant H2 contamination assessment for the Phase~1 sweep.}
\label{tab:variant-impact}
\end{table}

\paragraph{Warning.} All four variants run with \texttt{--damps\_apc 1} by
default (see the \texttt{Namespace} echoed in the first line of cell~16:
\texttt{damps\_apc=1}). Consequently:

\begin{quote}
The result $R@20 = 0.0851$ for the \texttt{combined} variant
\textbf{is not the true ceiling of the recipe}. It is the recipe \emph{plus}
phantom-APC noise. The real gap versus the paper's $0.0881$ could be smaller
\emph{or} larger than the reported $-3.4\%$; nothing definitive can be
concluded until either (i)~the true metadata file is restored, or
(ii)~APC is disabled entirely.
\end{quote}

\noindent This directly explains the stop-gate verdict produced by cell~20:
\textbf{all four variants FAIL Phase~1} with $R@20 < 0.0870$, which is
exactly the qualitative signature of H2 (hash-fallback weakens the cluster
signal).

\section{Relationship to the Roadmap's Top-3 Root Causes}

\begin{table}[h]
\centering
\renewcommand{\arraystretch}{1.25}
\begin{tabularx}{\textwidth}{|l|X|}
\hline
\textbf{Hypothesis} & \textbf{Status after this run} \\
\hline
\textbf{H2} -- APC phantom clusters via hash-fallback &
\textbf{STRONGLY CONFIRMED} (pre-flight banner + D1 exit~1) \\
\hline
\textbf{H3} -- \texttt{topk}=5 too small &
Inconclusive -- need the D2 sweep output for $K \in \{5, 10, 15, 20\}$ \\
\hline
\textbf{H10} -- sampled-eval vs.\ all-ranking mismatch &
Inconclusive -- need the D3 pure-MMHCL output (compare against the
$0.0881$ paper baseline) \\
\hline
\end{tabularx}
\caption{Updated status of the three highest-priority hypotheses.}
\label{tab:hypotheses}
\end{table}

\noindent The positive H2 verdict \textbf{raises the priority} of two actions:

\begin{enumerate}[leftmargin=*]
  \item Restore the original \texttt{meta\_categories.npy} (from the MMHCL
        paper release) and re-run the Phase~1 sweep. This is a
        \textbf{necessary precondition} for any further Phase~1 conclusion
        to be valid.
  \item Read the D2/D3 outputs that have already been computed (but not yet
        pasted) to decide whether $K$ needs to be increased \emph{after} the
        metadata is restored.
\end{enumerate}

\section{Why Did R@20 Still Reach $\sim$0.085 Under Hash-Fallback?}

This subtle point deserves a clear explanation. The hash-fallback in
\texttt{load\_data.py} generates clusters as
\(
\texttt{cluster}(i) = \texttt{hash}(\texttt{item\_id}) \;\bmod\; \texttt{num\_categories}.
\)
The resulting signal has structure but \emph{does not} correlate with the
true semantics of any item. When APC consumes it:

\begin{itemize}[leftmargin=*]
  \item the APC loss still converges -- the hash is deterministic, so the
        clusters remain stable across epochs;
  \item but the gradient signal propagated into item embeddings is
        \textbf{structured noise} -- the embeddings end up learning a
        ``fake feature'' whose effect on Recall@20 may help or hurt by
        chance.
\end{itemize}

The fact that \texttt{combined} ($0.0851$) is higher than \texttt{anchor}
($0.0795$) demonstrates that static-$\tau$ + AVRF-off is \emph{genuinely}
useful, because the hash-noise affects both variants symmetrically. But it
\textbf{cannot} be claimed that $0.0851$ is the ceiling of the recipe:
restoring the real metadata could push it above $0.087$ (passing the
stop-gate) or pull it down (if the paper's APC was tuned for the spurious
metadata in the first place).

This also suggests a useful \emph{secondary hypothesis}: if the real
metadata raises $R@20$ substantially, then APC is a genuinely valuable
module; if the real metadata barely changes anything, APC is mostly a
noise-regulariser and the module can be simplified or removed.

\section{Action Items in Priority Order}

\subsection*{1. IMMEDIATE -- Recover the D2/D3/APC outputs already computed
(do not waste 12~h of GPU)}

Open a Windows terminal in \texttt{MMHCL\_DAMPS\_Project/} and run:

\begin{lstlisting}
dir /s /b ..\Clothing\*diag_D2_topk*\*.txt
dir /s /b ..\Clothing\*diag_D3_mmhcl_pure*\*.txt
dir /s /b ..\Clothing\*diag_D1_apc_off*\*.txt
\end{lstlisting}

\noindent This should produce 6~\texttt{.txt} files
(4~D2 + 1~D3 + 1~APC-off). Share their contents, especially the lines:

\begin{lstlisting}
BEST_Test_Recall@20: ...
BEST_Test_NDCG@20:   ...
\end{lstlisting}

\noindent Alternatively, scroll back through the cell~15 output and copy
\textbf{everything} (not just the head and tail). Or, programmatically:

\begin{lstlisting}[style=py]
import json, nbformat
nb = nbformat.read(
    "Local_Random_seeds_train_mmhcl_clothing_colab_original.ipynb",
    as_version=4,
)
print(nb.cells[15].outputs[0]["text"])
\end{lstlisting}

\subsection*{2. RESTORE THE METADATA -- precondition for a valid Phase~1}

Two paths:

\begin{enumerate}[label=(\Alph*),leftmargin=*]
  \item \textbf{Download the original file from the MMHCL paper release.} The
        original MMHCL repository, distributed via the authors' Google Drive,
        typically ships \texttt{meta\_categories.npy} for Clothing, Sports
        and Baby. Place the file at:
        \begin{lstlisting}
data\Clothing\meta_categories.npy
        \end{lstlisting}
        Requirements: \texttt{len == 23033},
        \texttt{dtype} is integer, and at least 5 unique cluster labels.

  \item \textbf{Regenerate it from raw Amazon Clothing metadata.} If (A) is
        unavailable, the file can be regenerated by mapping each
        \texttt{item\_id} (range \(0\ldots 23032\)) to its top-level Amazon
        category
        (\texttt{Clothing, Shoes \& Jewelry > <subcategory>}).
        A \texttt{build\_meta\_categories.py} script can be provided on
        request.
\end{enumerate}

\subsection*{3. AFTER METADATA IS RESTORED -- re-run the Phase~1 sweep on
real clusters}

Re-execute the 20 runs of cell~16 (5 seeds $\times$ 4 variants $\times$ 250
epochs, $\approx 12$~h). Then compare:

\begin{itemize}[leftmargin=*]
  \item If \texttt{combined} reaches $R@20 \geq 0.0870$
        $\Rightarrow$ \textbf{stop-gate PASS}; Phase~1 closes and we move
        on to Phase~2.
  \item If it still FAILs but $R@20$ changes \emph{substantially}
        ($\Delta R@20 > 0.003$) $\Rightarrow$ APC carries a real signal;
        retune \texttt{damps\_warmup\_epochs} and
        \texttt{damps\_num\_categories}.
  \item If $R@20$ barely moves ($\Delta R@20 < 0.001$) $\Rightarrow$ APC
        is essentially noise regularisation; recommend dropping APC and
        moving directly to tuning $K$ or \texttt{rebuild\_R}.
\end{itemize}

\subsection*{4. OPTIONAL -- Update markdown cell~14 to record the finding}

Add the following note to cell~14:

\begin{quote}\itshape
The Phase~1 multi-seed sweep (cell~16) was executed with
\texttt{meta\_categories.npy} \textbf{MISSING} -- all 20 runs used the
hash-fallback APC clusters. The stop-gate verdict (cell~20: all variants
FAIL) is \textbf{confounded by H2} and must be re-evaluated after the
metadata file is restored.
\end{quote}

\section{Conclusion}

\begin{table}[h]
\centering
\renewcommand{\arraystretch}{1.3}
\begin{tabularx}{\textwidth}{|l|X|}
\hline
\textbf{Aspect} & \textbf{Assessment} \\
\hline
D1 / D2 / D3 infrastructure & Functioning to 100\% of spec \\
\hline
D1 verdict & H2 positive -- hash-fallback is active \\
\hline
Phase~1 sweep (cell~16) & Contaminated by H2 -- stop-gate verdict not final \\
\hline
D2 / D3 / APC outputs & Computed (\textasciitilde 12~h) but not yet shared -- must be recovered \\
\hline
Critical next step & Restore \texttt{meta\_categories.npy}; re-run the 20-run sweep \\
\hline
\end{tabularx}
\caption{One-line summary of the current state of Phase~1.}
\label{tab:summary}
\end{table}

\noindent Once the complete cell~15 output (or the contents of the 6 log
files under \verb|..\Clothing\diag_*|) is shared, definitive verdicts can be
issued for hypotheses \textbf{H3} and \textbf{H10}, and a precise execution
plan for Phase~1 round~2 can be drawn up.

\end{document}
